\chapter{Future Directions}
\label{ch:essay-future}

Future work divides into extensions that strengthen the applicability of existing theorems by providing additional empirical measurements and extensions that require new theoretical machinery beyond the admissibility framework developed here. The first category includes implementing projector and inverse-engine pairs for gesture, music, and audio modalities so that the cross-modal claims of Chapter~\ref{ch:essay-crossmodal} can transition from illustrative examples to measured results, running larger-scale versions of the Recoverability experiment to establish statistical confidence intervals for collision rates under different evolutionary conditions, and extending the Reconstruct experiment to test diagnostic predictions from Corollary~4.1 by deliberately constructing histories with known reachable-set relationships and measuring whether reconstruction uncertainty aligns with the theorem's predictions. The second category includes category-theoretic reformulation of the admissibility relation to make functorial properties explicit, investigation of hierarchical control structures in which higher-level reference conditions influence lower-level admissibility constraints as discussed in the perceptual-control literature, and extension of the framework to multi-scale systems where admissibility relations at one temporal or organizational scale constrain but do not fully determine admissibility at other scales. Planned experiments in the repository that would test the framework under progressively more challenging conditions include the False Cognate Laboratory, which measures accidental resemblance rates and tests evidential requirements against false-positive thresholds; the Dialect Continuum, which challenges tree-based genealogical assumptions by modeling diffusion through continuous space; Borrowing Without Ancestry, which introduces horizontal transmission to test when contact is mistaken for common descent; the Babel Experiment, which investigates whether linguistic differentiation emerges from initially uniform populations without encoding differentiation as an explicit outcome; Maximum Language, which complements the implemented Minimum Language by beginning with excessive complexity and measuring what survives imperfect transmission; Glyph Evolution, which extends the framework to writing-system evolution; ASL Handshape Drift, which tests whether the admissibility apparatus remains useful when linguistic evolution occurs in a different articulatory and perceptual modality; The Last Similar Thing, which compares recency-based transmission against retrieval based on structural and contextual similarity; and Language Earth, the long-term integration experiment combining migration, contact, divergence, borrowing, innovation, writing, prestige, isolation, and reconnection while retaining complete ground truth for subsequent reconstruction attempts. These experimental extensions would stress-test the distinction between inference error and genuine information loss under conditions including sparse evidence, convergent evolution, modality change, and non-tree genealogy, addressing not merely whether reconstruction succeeds but why it fails when it does. The framework's implications for other domains involving inverse problems under lossy projection may prove relevant to fields as distinct as paleontology, archaeology, and signal processing, though such applications remain speculative absent domain-specific validation of the admissibility structure's appropriateness.
